\subsection{Critérios de erro de atitude}
\label{sec:criterios_erro_atitude}

Os requisitos de apontamento da base são expressos em termos de três conjuntos de limites:

(i) erro angular máximo entre o referencial da base e o referencial desejado;
(ii) norma máxima da parte vetorial do quaternion de erro;
(iii) limites de velocidade angular em cada eixo do corpo.

O erro angular é medido como o menor ângulo de rotação que leva o referencial da base ao referencial desejado, devendo permanecer abaixo de um valor prescrito.
De forma equivalente, utiliza-se a parte vetorial do quaternion de erro, impondo-se um limite para sua norma.
Além disso, a velocidade angular da base deve permanecer dentro de limites máximos admissíveis em cada componente.
A atitude é considerada sincronizada sempre que todas essas grandezas permanecerem simultaneamente dentro dos limites especificados.
